%%%%%%%%%% source scan idx 78 | folio 71 %%%%%%%%%%

\underline{引理} 2.3. \underline{沿用 2.1 的记号, 假设 $X$ 由 $Q=0$ 定义.}

\begin{enumerate}[label=(\roman*),leftmargin=2.6em]
\item $\Omega_X^n\simeq\mathcal O_X(-n)$.

\item
\begin{equation*}
\operatorname{Pic}(X)\simeq
\begin{cases}
0, & n=0,\\
\mathbb Z, & n=1;\ \mathcal O_X(1)\text{ 是一个生成元的平方},\\
\mathbb Z\oplus\mathbb Z, & n=2,\\
\mathbb Z, & n\geq 3;\ \mathcal O_X(1)\text{ 是一个生成元}.
\end{cases}
\end{equation*}

\item 当 $i>0$ 时 $H^i(X,\mathcal O_X)=0$, 且 $H^1(X,\mathcal O_X(1))=0$.

\item 自然映射
\begin{equation*}
H^0\bigl(\mathbb P(V^*),\mathcal O_{\mathbb P(V^*)}(1)\bigr)
=V^*\xrightarrow{\sim}H^0(X,\mathcal O_X(1))
\end{equation*}
为同构; 且当 $n>0$ 时,
$H^0(X,\mathcal O_X)=k$.
\end{enumerate}

断言 (i) 由 $X$ 在 $\mathbb P(V^*)$ 中的余法层为 $\mathcal O_X(-2)$ 以及 $\mathbb P^r$ 的典范层为 $\mathcal O(-r-1)$ 得到. 断言 (ii) 在 $n\leq 2$ 时由 2.2 得到, 在 $n\geq 3$ 时由 SGA~2, XII, 3.7 得到. 最后, (iii) 和 (iv) 由 FAC, 第~78 节得到.

\underline{定义} 2.4. \underline{概形 $S$ 上维数为 $n$ 的光滑二次超曲面, 是指在 $S$ 上固有且光滑、其几何纤维均为光滑二次超曲面的 $S$-概形 $f:X\longrightarrow S$.}

当 $n=0$ 时, $S$ 上的零维光滑二次超曲面就是 $S$ 的二重平展覆盖. 当 $n=1$ 时, 它就是 $S$ 上一维的 Severi--Brauer 簇.

2.5. 设 $\rho:X\longrightarrow S$ 为 $S$ 上维数为 $n$ 的光滑二次超曲面. 由 2.3(i), 可逆层 $\Omega^n_{X/S}$ 的逆为丰沛层.

%%%%%%%%%% source scan idx 79 | folio 72 %%%%%%%%%%

由 2.3(iii), Picard 概形 $\operatorname{Pic}_{X/S}$ 在 $S$ 上光滑 (因为 $H^2(X,\mathcal O_X)=0$)、在 $S$ 上非分歧 (因为 $H^1(X,\mathcal O_X)=0$), 且在 $S$ 上本质固有 (因为 $X/S$ 光滑). 因此, 在 $S$ 的平展拓扑下局部地, 它是具有 2.3(ii) 所述结构的常离散群.\footnote{印本引用 2.3(ii) 来说明上同调消失, 引用 2.3(iii) 来计算 Picard 群. 这里对调了这两个引用.}

特别地, 若 $X$ 有截面, 则 $X$ 上存在可逆层 $L$, 使得 $L^{\otimes(-n)}$ 与 $\Omega^n_{X/S}$ 在 $\operatorname{Pic}_{X/S}$ 中的像相同; 任意两个这样的层在 $S$ 上局部同构. 由 2.3(iii) 和 [1], 0.5, 第~19 页, $\rho_*L$ 局部自由, 且其构造与每个基变换交换. 因而 $\rho_*L$ 局部自由、秩为 $n+2$, 映射 $\rho^*\rho_*L\longrightarrow L$ 为满射, 且典范映射
\begin{equation*}
X\longrightarrow\mathbb P(\rho_*L)
\end{equation*}
为闭浸入. 它将 $X$ 等同于 $\mathbb P(\rho_*L)$ 中次数为二的相对除子; 在 $S$ 上局部地, $X$ 是由 $(\rho_*L)^*$ 上的普通二次型定义的 $\mathbb P(\rho_*L)$ 中的二次超曲面, 此型在相差可逆因子的意义下唯一.

\footnote{印本在以 $\rho$ 表示投影之后, 又在直像公式中将其符号改为 $p$. 此处统一了记号; 对 2.3(iii) 的引用则予以保留.}

2.6. 假设 $n>0$. 若 $L_1$ 和 $L_2$ 为上述类型的两个可逆层, 则在 $S$ 上局部地存在同构 $u:L_1\xrightarrow{\sim}L_2$. 由 2.3(iv) 知 $\rho_*\mathcal O_X=\mathcal O_S$, 所以该同构在相差可逆因子的意义下唯一. 因此所有同构 $u:L_1\xrightarrow{\sim}L_2$ 都诱导相同的同构
\begin{equation*}
\mathbb P(u):\mathbb P(\rho_*L_1)\xrightarrow{\sim}
\mathbb P(\rho_*L_2).
\end{equation*}
于是 $\mathbb P(\rho_*L_1)$ 与 $\mathbb P(\rho_*L_2)$ 典范同构, 且只依赖于 $X/S$, 与 $L$ 的选择无关. 我们将其记为 $P(X)$.

%%%%%%%%%% source scan idx 80 | folio 73 %%%%%%%%%%

射影空间 $P(X)$ 在 $S$ 的平展拓扑下局部地有定义. 由下降, 它定义了 $S$ 上一个仍记为 $P(X)$ 的 Severi--Brauer 簇, 其中 $X$ 为次数为二的相对除子.

在 $S$ 上局部地, 对 $X\hookrightarrow P(X)$ 定义了 Clifford 代数 $C^+(X,P(X))$ (1.3). 通过下降得到 $S$ 上的代数, 我们将其简记为 $C^+(X)$.

2.7. 假设 $n=2m>0$. 则 $C^+(X)$ 的中心 $Z$ 为 $\mathcal O_S$-代数, 它定义了 $S$ 的二重平展覆盖 $Z(X)$.

$X$ 的一个\underline{母空间}(generator)是指 $X$ 的子概形 $D$, 它是 $P(X)$ 中维数为 $m$ 的线性子空间.

在 $S$ 上局部地, 当 $P(X)\simeq\mathbb P(V^*)$ 且 $X$ 由 $Q=0$ 定义时, 按定义, $X$ 的母空间 $D$ 等同于 $V$ 中局部为直和项、全奇异且秩为 $m+1$ 的子模 $W(D)$. 由 1.8, $C^+(X)$ 中心内的幂等元 $e(W(D))$ 只依赖于 $D\subset X$. 以 $e(D)$ 表示 $Z(X)$ 的相应截面, 即 $e(W(D))$ 在其上取值为一的截面.

$X/S$ 的母空间是某个 $S$ 上射影概形 $\operatorname{Gen}(X)$ 的 $S$-点, 而 $e$ 定义 $S$-概形的态射

\begin{equation}\tag{2.7.1}
e:\operatorname{Gen}(X)\longrightarrow Z(X).
\end{equation}

当 $n=0$ 时, 置 $Z(X)=X$, 将 $X$ 的截面称为 $X$ 的母空间, 并将 $e$ 定义为恒等映射.

%%%%%%%%%% source scan idx 81 | folio 74 %%%%%%%%%%

\underline{命题} 2.8. \underline{在 2.7 的假设下, $p:\operatorname{Gen}(X)\longrightarrow S$ 为射影且光滑的态射, 而 (2.7.1) 为其 Stein 分解.}

这一问题对于 $S$ 的平展拓扑是局部的. 局部地, 若 $D$ 为一个母空间, 则由 1.6, 与 $D$ 不相交的母空间 $D'$ 可等同于适当的 $S$ 上仿射空间的 $S$-点. 由此, $\operatorname{Gen}(X)$ 被在 $S$ 上光滑的开子概形覆盖.

为证明 $e$ 的几何纤维几何连通, 只须考虑 $S$ 为代数闭域 $k$ 的谱, 且 $X$ 由 $\mathbb P(V^*)$ 中的 $Q=0$ 定义的情形.

设 $W_1$ 和 $W_2$ 为 $V$ 中秩为 $m+1$ 的两个全迷向子空间, 满足 $e(W_1)=e(W_2)$, 即 $W_1/(W_1\cap W_2)$ 的维数为偶数 (1.12). 与 1.12 中相同, 得到正交分解
\begin{equation*}
V=V_0\oplus\bigoplus_i V_i
\end{equation*}
使得
\begin{enumerate}[label=\alph*),leftmargin=2.5em]
\item 对 $\alpha=1,2$, 有 $W_\alpha=(W_\alpha\cap V_0)\oplus
  \displaystyle\bigoplus_i(W_\alpha\cap V_i)$;
\item $W_1\cap V_0=W_2\cap V_0$;
\item $\dim V_i=4$ 且
  $V_i=(W_1\cap V_i)\oplus(W_2\cap V_i)$.
\end{enumerate}
只须再证明 $W_1\cap V_i$ 与 $W_2\cap V_i$ 属于 $V_i$ 的极大全奇异子空间的同一个连通族. 换言之, 断言归结为二维二次超曲面的情形. 这样的二次超曲面同构于 $\mathbb P^1\times\mathbb P^1$, 其母空间为子概形 $\{t\}\times\mathbb P^1$ 和 $\mathbb P^1\times\{t\}$; 此时命题显然成立.
